\section{Related Work}
\label{sec:related}

\subsection{From Mental-State Inference to Functional Use}

Machine-ToM benchmarks have made latent-state inference increasingly controlled and difficult. ToMi exposes shortcuts in false-belief question answering \citep{le2019revisiting}; BigToM uses causal templates \citep{gandhi2023bigtom}; and FANToM tests consistency across questions that depend on the same information asymmetry \citep{kim2023fantom}. Hi-ToM, OpenToM, and ToMBench broaden reasoning order, narrative variation, and the taxonomy of social-cognitive abilities \citep{wu2023hitom,xu2024opentom,chen2024tombench}.

These advances strengthen the validity of inference from \emph{supplied} evidence. They do not make evidence acquisition part of the tested policy. SimpleToM identifies a related boundary: explicitly inferring a mental state does not guarantee that a model will use it when predicting or judging behavior \citep{gu2024simpletom}. \knowact takes this boundary as a design requirement. A user model must affect the next diagnostic action, and the result of that process must remain directly scoreable.

Interactive social benchmarks test behavior in richer loops. SOTOPIA and AgentSense evaluate social goals, strategic communication, and private-information reasoning in multi-turn scenarios \citep{zhou2024sotopia,mou2025agentsense}. Applied-ToM agents explicitly insert mental-state hypotheses into social, persuasive, or software-engineering policies \citep{hwang2026infusing,wang2026tomelp,zhou2026tomswe}.

Their question is whether user modeling improves an external outcome. Ours is whether the user model itself is acquired accurately and adaptively. These are complementary claims: downstream success establishes usefulness, whereas direct reconstruction establishes what was learned about the user. Broader agent evaluation likewise shows that single-turn ability and final success can obscure multi-turn and trajectory-level differences \citep{wang2024mint,ma2024agentboard}.

\subsection{Passive User Modeling and Active Diagnosis}

Personalization benchmarks infer preferences and traits from user histories. LaMP evaluates whether profile information improves personalized outputs \citep{salemi2024lamp}; PERSONAMEM asks whether models track profiles that evolve across long interaction histories \citep{jiang2025personamem}. In both cases, the history is evidence to be interpreted. The tested model does not decide which observation would most reduce uncertainty about the user.

Knowledge tracing provides the closest formal analogue. Bayesian and neural tracing methods estimate latent mastery from sequences of item responses \citep{corbett1994knowledgetracing,piech2015dkt,ghosh2020akt}. LongTutor further separates historical evidence acquisition, knowledge diagnosis, and adaptive teaching, revealing that strength at one stage need not transfer to the next \citep{li2026longtutor}.

\knowact changes who controls the evidence. The tested agent asks open-ended questions and is scored on the full latent state after a bounded number of observations. The benchmark is therefore an active diagnosis problem rather than next-response prediction over a fixed curriculum.

\subsection{LLM-Based User Simulation}

User simulation trades experimental control for behavioral fidelity. Agenda-based simulators are reproducible but narrow; LLM simulators are expressive but can invent goals, leak hidden state, or impose model-specific artifacts. DuetSim separates generation and verification to improve response quality \citep{luo2024duetsim}. SimulatorArena instead evaluates the proxy assumption directly, comparing simulated behavior and assistant rankings with human interaction \citep{dou2025simulatorarena}. Earlier simulator research makes the same conceptual point: fluency is not fidelity to task-relevant behavior \citep{shi2019usersimulators}.

\knowact does not treat an LLM simulator as ground truth about human conversation. It uses simulation to hold a reviewed latent state fixed while agents choose evidence. Structural access controls limit what can be expressed, deterministic scoring prevents the simulator from judging success, and matched human episodes test whether conclusions transfer beyond the proxy. Simulator validation is therefore part of the benchmark claim, not an optional naturalness check.

\begin{table}[t]
\caption{Positioning of \knowact relative to adjacent evaluation paradigms. ``Active probes'' means the tested model chooses what evidence to request.}
\label{tab:positioning}
\centering
\small
\begin{tabularx}{\linewidth}{>{\raggedright\arraybackslash}Xcccc}
\toprule
\textbf{Paradigm} & \textbf{Interactive} & \shortstack{\textbf{Explicit latent}\\\textbf{user state}} & \shortstack{\textbf{Active}\\\textbf{probes}} & \shortstack{\textbf{Direct state}\\\textbf{score}} \\
\midrule
Narrative / conversational ToM & Limited & Partial & No & Usually QA \\
Social-agent arenas & Yes & Usually no & Yes & No \\
Personalization from histories & Usually no & Profile / preference & No & Sometimes \\
Knowledge tracing & Sequential & Mastery & Usually no & Prediction-based \\
\textbf{\knowact} & \textbf{Yes} & \textbf{Knowledge map} & \textbf{Yes} & \textbf{Yes} \\
\bottomrule
\end{tabularx}
\end{table}
